\subsection{核型映射}

在本节中，F 表示 K 上的希尔伯特空间。当 K = C 时，我们回顾（I, p. 139，定义 3），对 \(u \in \mathcal{L}(\mathrm{E}; \mathrm{F})\) ，用 \(|u|\) 表示 E 的正自同态 \(\sqrt{u^{*} \circ u}\) 。当 K = R 时，元素 \(|u_{(\mathbf{C})}|\) （\(\mathcal{L}(\mathrm{E}_{(\mathbf{C})})\) 中的元素）具有形式 \(v_{(\mathbf{C})}\) ，其中自同态 \(v \in \mathcal{L}(\mathrm{E})\) 是唯一的且仍记作 \(|u|\) （参见 I, p. 87）。

我们用 \((u,v)\mapsto\langle u|v\rangle=\mathrm{Tr}(u^{*}v)\) 表示希尔伯特空间 \(\mathcal{L}_{2}(\mathrm{E};\mathrm{F})\) 中的内积（EVT, V, p. 53，注 1 与注 2）。

对每个 \(u \in \mathcal{L}(\mathrm{E};\mathrm{F})\) ，用 \(\| u\| _1\) 表示量 \(\mathrm{Tr}(|u|)\) （若 \(|u|\) 有有限迹）；否则置 \(\| u\| _1 = +\infty\) 。于是我们有 \(\| u\| _1 \in \mathbf{R}_+ \cup \{+\infty\}\) 。由于 \(\| u\| = \| |u|\|\) （I, p. 139 的命题 10, a)），且当 \(v\) 是正的并有限迹时有 \(\mathrm{Tr}(v) \geqslant \| v\|\) （EVT, V, p. 49，公式 (24bis) 及 p. 44，命题 9），由此可得

\[
\| u \| \leqslant \| u \| _ {1}.\tag{9}
\]

若 \(\| u\| _1\) 有限，则 \(u\) 是希尔伯特--施密特算子，且 \(\| u\| _2 \leqslant \| u\| _1\) （IV, p. 165 的推论 4）。

命题 14. --- 设 \(u \in \mathcal{L}_2(\mathrm{E};\mathrm{F})\) 。我们有

\[
\| u\|_{1} = \sup_{\substack{v\in \mathcal{L}_{2}(\mathrm{E};\mathrm{F})\\ \| v\| \leqslant 1}}|\langle v  |  u\rangle |,
\]

其中的上确界在 \(\mathbf{R}_{+} \cup \{+\infty\}\) 中计算。

设 \(\mathrm{B} = (e_i)_{i \in \mathrm{I}}\) 是 u 的始空间 \(\operatorname{Ker}(u)^{\circ}\) 的标准正交基，\(\mathrm{C} = (f_i)_{i \in \mathrm{I}}\) 是 F 中的标准正交族，\((\alpha_i)_{i \in \mathrm{I}}\) 是 \((\mathbf{R}_+^*)^{\mathrm{I}}\) 中的族，使得 \(u(e_i) = \alpha_i f_i\) 对所有 \(i \in I\) 成立（IV, p. 150 的推论 2）。族 \((\alpha_i)_{i \in \mathrm{I}}\) 平方可积，因为 u 属于 \(\mathcal{L}_2(\mathrm{E};\mathrm{F})\) （IV, p. 165 的推论 4）。设 \((e_j)_{j \in \mathrm{J}}\) 是 \((e_i)_{i \in \mathrm{I}}\) 的扩张，即 E 的标准正交基。

设 L 是 I 的有限子集。设 \(v_{L}\) 是由

\[
v _ {\mathrm{L}} (x) = \sum_ {i \in \mathrm{L}} \langle e _ {i} | x \rangle f _ {i}
\]

对所有 \(x\) 属于 E 定义的从 E 到 F 中的连续有限秩线性映射。我们有 \(\| v_{\mathrm{L}}\| \leqslant 1\) 且 \(v_{\mathrm{L}}\in \mathcal{L}_{2}(\mathrm{E};\mathrm{F})\) 。此外，对每个 \(j\in \mathrm{J}\) ，有 \(v_{\mathrm{L}}(e_j) = 0\) （若 \(j\notin \mathrm{L}\) ），且有 \(v_{\mathrm{L}}(e_j) = f_j\) （若 \(j\in \mathrm{L}\) ）。于是

\[
| \langle v _ {\mathrm{L}} | u \rangle | = | \operatorname{Tr} (v _ {\mathrm{L}} ^ {*} u) | = \left| \sum_ {j \in \mathrm{J}} \langle v _ {\mathrm{L}} (e _ {j}) | u (e _ {j}) \rangle \right| = \sum_ {j \in \mathrm{L}} \alpha_ {j}.
\]

由同前，我们推出

\[
\| u\|_{1} = \sup_{\text{L}}\sum_{j\in \text{L}}\alpha_{j}\leqslant \sup_{\substack{v\in \mathcal{L}_{2}(\text{E};\text{F})\\ \| v\| \leqslant 1}}|\langle v\mid u\rangle |\tag{10}
\]

在 \(\mathbf{R}_{+}\cup\{+\infty\}\) 中成立。

当 \(\|u\|_{1}=+\infty\) 时，这即给出命题中的等式。

设 \(\|u\|_{1}\) 有限。对每个 \(i \in I\) ，设 \(p_{i}\) 是从 E 到 F 的线性映射，满足 \(p_{i}(e_{i}) = f_{i}\) 且 \(p_{i}(x) = 0\) 对所有 \(x \in e_{i}^{\circ}\) 成立。对 I 中所有的 j 和 k，我们有

\[
\langle p _ {j} \mid p _ {k} \rangle = \mathrm{Tr} (p _ {j} ^ {*} p _ {k}) = \sum_ {i \in \mathrm{J}} \langle p _ {j} (e _ {i}) \mid p _ {k} (e _ {i}) \rangle ,
\]

除非 \(j = k\) ，它为零；在此情形这个量等于 \(\| f_k\|^2 = 1\) 。因此族 \((p_i)_{i \in \mathrm{I}}\) 在 \(\mathcal{L}_2(\mathrm{E};\mathrm{F})\) 中标准正交。从而族 \((\alpha_i p_i)_{i \in \mathrm{I}}\) 在 \(\mathcal{L}_2(\mathrm{E};\mathrm{F})\) 中可和；其和等于 \(u\) ，因为这两个连续线性映射在元素 \(e_i\) 上（对所有 \(i \in \mathrm{J}\)）相合。

设 \(v \in \mathcal{L}_2(\mathrm{E};\mathrm{F})\) 。对每个 \(i \in \mathrm{I}\) ，由此可得

\[
\langle v \mid p _ {i} \rangle = \operatorname{Tr} (v ^ {*} p _ {i}) = \sum_ {j \in \mathrm{J}} \langle v (e _ {j}) \mid p _ {i} (e _ {j}) \rangle = \langle v (e _ {i}) \mid f _ {i} \rangle .
\]

若 \(\|v\| \leqslant 1\) ，我们有

\[
\begin{array}{c} | \langle v | u \rangle | = \left| \langle v | \sum_ {i \in I} \alpha_ {i} p _ {i} \rangle \right| \leqslant \sum_ {i \in I} \alpha_ {i} | \langle v | p _ {i} \rangle | \\ = \sum_ {i \in I} \alpha_ {i} | \langle v (e _ {i}) | f _ {i} \rangle | \leqslant \sum_ {i \in I} \alpha_ {i} = \mathrm{Tr} (| u |), \end{array}
\]

由此得到不等式

\[
\sup_{\substack{v\in \mathcal{L}_{2}(\mathrm{E};\mathrm{F})\\ \| v\| \leqslant 1}}|\langle v\mid u\rangle |\leqslant \| u\|_{1},
\]

它与公式 (10) 相结合，便完成了命题的证明。

由这个命题可得，集 \(\mathcal{L}_1(\mathrm{E};\mathrm{F})\) （从 E 到 F 的连续线性映射 \(u\) 中使 \(\| u\| _1\) 有限者所成之集）是 \(\mathcal{L}_{2}(\mathrm{E};\mathrm{F})\) 的向量子空间，且映射 \(u\mapsto \| u\| _1\) 是 \(\mathcal{L}_1(\mathrm{E};\mathrm{F})\) 上的半范数；不等式 (9) 表明它是一个范数。

定义 4. --- 我们称向量空间 \(\mathcal{L}_{1}(\mathrm{E};\mathrm{F})\) （赋以范数 \(u\mapsto \| u\|_{1}\)）为从 E 到 F 的核型映射所成的空间。若 \(u\in \mathcal{L}_1(\mathrm{E};\mathrm{F})\) ，则称 \(u\) 是核型的。

注记. --- 1) 假设 K = R。设 \(u \in \mathcal{L}(E; F)\) 。我们有 \(u \in \mathcal{L}_{1}(E; F)\) 当且仅当 \(u_{(\mathbf{C})} \in \mathcal{L}_{1}(E_{(\mathbf{C})}; F_{(\mathbf{C})})\) ；此时有 \(\|u\|_{1} = \|u_{(\mathbf{C})}\|_{1}\) 。

\begin{enumerate}
\def\labelenumi{\arabic{enumi})}
\setcounter{enumi}{1}
\tightlist
\item
  设 \(u\) 是从 E 到 F 的核型映射。由于映射 \(u\) 是希尔伯特--施密特的，它是紧的（IV, p. 165 的推论 2）。此外，命题 14 蕴含对每个 \(v \in \mathcal{L}_2(\mathrm{E};\mathrm{F})\) ，我们有
\end{enumerate}

\[
| \langle v | u \rangle | \leqslant \| v \| \| u \| _ {1}.\tag{11}
\]

\begin{enumerate}
\def\labelenumi{\arabic{enumi})}
\setcounter{enumi}{2}
\item
  假设 \(\mathrm{K} = \mathbf{C}\)。设 \(u \in \mathcal{L}_1(\mathrm{E};\mathrm{E})\) 且 \(u\) 是正的。范数 \(\| u\| _1\) （\(u\) 的范数）是序列 \((\lambda_n(u))_{n\in \mathrm{I_E}}\) 的和（IV, p. 165 的推论 3）。
\item
  \(\mathcal{L}_1(\mathrm{E};\mathrm{F})\) 到 \(\mathcal{L}(\mathrm{E};\mathrm{F})\) 中的典范单射的范数 \(\leqslant 1\) （p. 167 的不等式 (9)）。
\end{enumerate}

命题 15. --- 设 G 是 K 上的希尔伯特空间。映射 \((u,v)\mapsto v\circ u\) （从 \(\mathcal{L}(\mathrm{E};\mathrm{F})\times\mathcal{L}(\mathrm{F};\mathrm{G})\) 到 \(\mathcal{L}(\mathrm{E};\mathrm{G})\) 中的映射）过渡到商，定义了范数 \(\leqslant1\) 的从 \(\mathcal{L}_{2}(\mathrm{E};\mathrm{F})\times\mathcal{L}_{2}(\mathrm{F};\mathrm{G})\) 到 \(\mathcal{L}_{1}(\mathrm{E};\mathrm{G})\) 中的连续双线性映射。

设 \(u \in \mathcal{L}_2(\mathrm{E};\mathrm{F})\) 且 \(v \in \mathcal{L}_2(\mathrm{F};\mathrm{G})\)。设 \(w \in \mathcal{L}_2(\mathrm{E};\mathrm{G})\) 。我们得到 \(\langle uv | w \rangle = \operatorname{Tr}(v^* u^* w) = \langle v | u^* w \rangle\) ，由此

\[
| \langle u v | w \rangle | \leqslant \| v \| _ {2} \| u ^ {*} w \| _ {2} \leqslant \| v \| _ {2} \| u \| _ {2} \| w \|
\]

此处用到柯西--施瓦茨不等式及 EVT, V, p. 52 的公式 (37)。因此结论由命题 14 得出。

引理 8. --- 映射 \(u \mapsto u^{*}\) （从 \(\mathcal{L}(\mathrm{E};\mathrm{F})\) 到 \(\mathcal{L}(\mathrm{F};\mathrm{E})\) 中的映射）过渡到商，定义了从 \(\mathcal{L}_{1}(\mathrm{E};\mathrm{F})\) 到 \(\mathcal{L}_{1}(\mathrm{F};\mathrm{E})\) 中的等距线性映射。

设 \(u \in \mathcal{L}_1(\mathrm{E};\mathrm{F})\) 。由于 \(u\) 是希尔伯特--施密特映射，\(u^*\) 也是（EVT, V, p. 54）。映射 \(v \mapsto v^*\) 是从 \(v \in \mathcal{L}_2(\mathrm{E};\mathrm{F})\) 中满足 \(\|v\| \leqslant 1\) 者所成之集到 \(w \in \mathcal{L}_2(\mathrm{F};\mathrm{E})\) 中满足 \(\|w\| \leqslant 1\) 者所成之集上的双射；对每个 \(v \in \mathcal{L}_2(\mathrm{E};\mathrm{F})\)

满足 \(\|v\|\leqslant1\) 者，我们有 \(\langle v|u\rangle=\langle u^{*}|v^{*}\rangle\) （EVT, V, p. 54，公式 (42)），于是结论由命题 14 得出。

命题 16. --- 设 \(E_{1}\) 与 \(F_{1}\) 是希尔伯特空间。设 \(u\) 属于 \(\mathcal{L}_{1}(\mathrm{E};\mathrm{F})\) ，\(v\) 属于 \(\mathcal{L}(\mathrm{E}_{1};\mathrm{E})\) ，\(v_{2}\) 属于 \(\mathcal{L}(\mathrm{F};\mathrm{F}_{1})\)。则有 \(v_{2}uv_{1} \in \mathcal{L}_{1}(\mathrm{E}_{1};\mathrm{F}_{1})\) 且 \(\|v_{2}uv_{1}\|_{1} \leqslant \|v_{2}\|\|v_{1}\|\|u\|_{1}\) 。

设 \(w \in \mathcal{L}_2(\mathrm{E};\mathrm{F}_1)\) 满足 \(\| w\| \leqslant 1\) 。我们有 \(v_{2}u \in \mathcal{L}_{2}(\mathrm{E};\mathrm{F}_{1})\) （EVT, V, p. 52，公式 (36)）。由于 \(v_{2}^{*}w \in \mathcal{L}_{2}(\mathrm{E};\mathrm{F})\) （同前），由此可得

\[
| \langle w | v _ {2} u \rangle | = | \langle v _ {2} ^ {*} w | u \rangle | \leqslant \| v _ {2} \| \| u \| _ {1}
\]

（公式 (11)），于是 \(v_2u \in \mathcal{L}_1(\mathrm{E};\mathrm{F}_1)\) 且 \(\| v_2u\| _1 \leqslant \| v_2\| \| u\| _1\) （命题 14）。设 \(v_{1} \in \mathcal{L}(\mathrm{E}_{1};\mathrm{E})\) ；由于 \(uv_{1} = (v_{1}^{*}u^{*})^{*}\) ，我们有 \(uv_{1} \in \mathcal{L}_{1}(\mathrm{E}_{1};\mathrm{F})\) 且 \(\| uv_{1}\|_{1} \leqslant \| v_{1}^{*}\| \| u^{*}\|_{1} = \| v_{1}\| \| u\|_{1}\) （引理 8）。

命题由这些不等式立即得出。

命题 17. — 空间 \(\mathcal{L}_{1}(\mathrm{E};\mathrm{E})\) 与 E 的有限迹自同态所成的空间 \(\mathcal{L}_{1}(\mathrm{E})\) 相重合，且我们有 \(|\operatorname{Tr}(u)|\leqslant \| u\|_{1}\) ，对 E 的每个有限迹自同态 \(u\) 成立。

我们可以假设 K = C。按定义，空间 \(\mathcal{L}_{1}(E;E)\) 包含有限迹正自同态所成之集，因此 \(\mathcal{L}_{1}(E)\subset\mathcal{L}_{1}(E;E)\) （EVT, p. 50，定义 8）。反过来，我们要证明 \(\mathcal{L}_{1}(E;E)\) 含于 \(\mathcal{L}_{1}(E)\) 。

设 \(u \in \mathcal{L}_{1}(\mathrm{E};\mathrm{E})\) 。我们有 \(u^{*} \in \mathcal{L}_{1}(\mathrm{E};\mathrm{E})\) （引理 8）；因此只需证明 \(\mathcal{L}_{1}(\mathrm{E};\mathrm{E})\) 中的埃尔米特元素有有限迹（I, p. 96 的引理 2）。

设 u 是这样一个自同态。它是紧的（p. 169 的注 2）。设 B 是 E 的标准正交基，使得 u 在基 B 中是对角的（IV, p. 149 的定理 1），并设 \(\lambda\) 是 u 关于 B 的特征值族。自同态 \(|u|\) 在基 B 中可对角化，其特征值族是 \(|\lambda|\) （IV, p. 147 的命题 1, d)）。由于 \(|u|\) 有有限迹，后一族可和（IV, p. 164 的命题 12），因此族 \(\lambda\) 可和。于是 u 有有限迹（同前），且我们有 \(|\mathrm{Tr}(u)| \leqslant \mathrm{Tr}(|u|) = \|u\|_{1}\) 。

从而，空间 \(\mathcal{L}_{1}(\mathrm{E})\) 是 C*-代数 \(\mathcal{L}(\mathrm{E})\) 的自伴双边理想。若 E 是无限维的，这个理想在 \(\mathcal{L}(\mathrm{E})\) 中不闭。

命题 17 与命题 16 表明，映射 \(b\) ——从 \(\mathcal{L}_1(\mathrm{E};\mathrm{F})\times \mathcal{L}(\mathrm{F};\mathrm{E})\) 到 \(\mathbf{C}\) 中的映射——由 \(b(u,v) = \operatorname {Tr}(vu)\) 定义，是连续双线性型，它使空间 \(\mathcal{L}_1(\mathrm{E};\mathrm{F})\) 与 \(\mathcal{L}(\mathrm{F};\mathrm{E})\) 处于对偶。

引理 9. --- 设 \(u \in \mathcal{L}_1(\mathrm{E};\mathrm{F})\) 。我们有

\[
\| u\|_{1} = \sup_{\substack{v\in \mathcal{L}^{c}(\mathrm{F};\mathrm{E})\\ \| v\| \leqslant 1}}|b(u,v)|.
\]

由于每个希尔伯特--施密特映射都是紧的（IV, p. 165 的推论 2），我们有

\[
\| u\|_{1}\leqslant \sup_{\substack{v\in \mathcal{L}^{\mathrm{c}}(\mathrm{E};\mathrm{F})\\ \| v\| \leqslant 1}}|\operatorname{Tr}(vu)| = \sup_{\substack{v\in \mathcal{L}^{\mathrm{c}}(\mathrm{F};\mathrm{E})\\ \| v\| \leqslant 1}}|b(u,v)|
\]

（命题 14）。

设 \(v \in \mathcal{L}^{c}(\mathrm{E};\mathrm{F})\) 的范数 \(\leqslant 1\) 。由于 \(\mathcal{L}_{2}(\mathrm{E};\mathrm{F})\) 在 \(\mathcal{L}^{c}(\mathrm{E};\mathrm{F})\) 中稠密，存在序列 \((v_{n})_{n\in\mathbf{N}}\) （\(\mathcal{L}_{2}(\mathrm{E};\mathrm{F})\) 中的序列）在 \(\mathcal{L}(\mathrm{E};\mathrm{F})\) 中收敛到 v。序列 \((b(u,v_{n}))_{n\in\mathbf{N}}\) 收敛到 \(b(u,v)\) ；由于 \(|b(u,v_{n})| = |\operatorname{Tr}(v_{n}u)| = |\langle v_{n}^{*}|u\rangle| \leqslant \|u\|_{1}\) （命题 14）对所有 \(n \in N\) 成立，由此可得 \(|b(u,v)| \leqslant \|u\|_{1}\) 。

用 \(\theta\) 表示连续线性映射

\[
\theta \colon \mathcal {L} _ {1} (\mathrm{E}; \mathrm{F}) \to \mathcal {L} ^ {\mathrm{c}} (\mathrm{F}; \mathrm{E}) ^ {\prime}
\]

使得 \(\theta(u)(v) = b(u,v)\) 。

命题 18. --- 映射 θ 是等距同构。

由引理 9，映射 \(\theta\) 是等距的，只需证明 \(\theta\) 是满射。

设 \(\lambda \in \mathcal{L}^{\mathrm{c}}(\mathrm{F};\mathrm{E})'\) 。由于 \(\| v\| \leqslant \| v\|_{2}\) 对所有 \(v\in \mathcal{L}_2(\mathrm{F};\mathrm{E})\) 成立（EVT, V, p. 52，公式 (33)），\(\lambda\) 在子空间 \(\mathcal{L}_2(\mathrm{F};\mathrm{E})\) 上的限制是希尔伯特空间 \(\mathcal{L}_2(\mathrm{F};\mathrm{E})\) 上的连续线性型。因此存在元素 \(u\) （\(\mathcal{L}_2(\mathrm{F};\mathrm{E})\) 的元素），使得 \(\lambda (v) = \langle u|v\rangle\) 对所有 \(v\in \mathcal{L}_2(\mathrm{F};\mathrm{E})\) 成立（EVT, V, p. 15，定理 3）。

对每个 \(w \in \mathcal{L}_2(\mathrm{E};\mathrm{F})\) （范数 \(\leqslant 1\) 者），我们有 \(|\langle w|u\rangle| = |\lambda(w)| \leqslant \| \lambda \|\) 。于是 \(u\) 是从 \(\mathrm{E}\) 到 \(\mathrm{F}\) 中的核型映射（命题 14）。

连续线性型 \(\lambda\) 与 \(\theta(u)\) 在稠密子空间 \(\mathcal{L}_{2}(\mathrm{F},\mathrm{E})\) （\(\mathcal{L}^{\mathrm{c}}(\mathrm{F};\mathrm{E})\) 的稠密子空间）上相等。由此可得 \(\lambda=\theta(u)\) ，证明完成。

推论. --- 赋范空间 \(\mathcal{L}_{1}(\mathrm{E};\mathrm{F})\) 是巴拿赫空间。

由于空间 \(\mathcal{L}^{\mathrm{c}}(\mathrm{F};\mathrm{E})\) 是赋范空间，断言由本命题及 EVT, III, p. 24 的推论 2 得出。

